% Options for packages loaded elsewhere
\PassOptionsToPackage{unicode}{hyperref}
\PassOptionsToPackage{hyphens}{url}
\documentclass[
]{article}
\usepackage{xcolor}
\usepackage[margin=1in]{geometry}
\usepackage{amsmath,amssymb}
\setcounter{secnumdepth}{5}
\usepackage{iftex}
\ifPDFTeX
  \usepackage[T1]{fontenc}
  \usepackage[utf8]{inputenc}
  \usepackage{textcomp} % provide euro and other symbols
\else % if luatex or xetex
  \usepackage{unicode-math} % this also loads fontspec
  \defaultfontfeatures{Scale=MatchLowercase}
  \defaultfontfeatures[\rmfamily]{Ligatures=TeX,Scale=1}
\fi
\usepackage{lmodern}
\ifPDFTeX\else
  % xetex/luatex font selection
\fi
% Use upquote if available, for straight quotes in verbatim environments
\IfFileExists{upquote.sty}{\usepackage{upquote}}{}
\IfFileExists{microtype.sty}{% use microtype if available
  \usepackage[]{microtype}
  \UseMicrotypeSet[protrusion]{basicmath} % disable protrusion for tt fonts
}{}
\makeatletter
\@ifundefined{KOMAClassName}{% if non-KOMA class
  \IfFileExists{parskip.sty}{%
    \usepackage{parskip}
  }{% else
    \setlength{\parindent}{0pt}
    \setlength{\parskip}{6pt plus 2pt minus 1pt}}
}{% if KOMA class
  \KOMAoptions{parskip=half}}
\makeatother
\usepackage{longtable,booktabs,array}
\newcounter{none} % for unnumbered tables
\usepackage{calc} % for calculating minipage widths
% Correct order of tables after \paragraph or \subparagraph
\usepackage{etoolbox}
\makeatletter
\patchcmd\longtable{\par}{\if@noskipsec\mbox{}\fi\par}{}{}
\makeatother
% Allow footnotes in longtable head/foot
\IfFileExists{footnotehyper.sty}{\usepackage{footnotehyper}}{\usepackage{footnote}}
\makesavenoteenv{longtable}
\usepackage{graphicx}
\makeatletter
\newsavebox\pandoc@box
\newcommand*\pandocbounded[1]{% scales image to fit in text height/width
  \sbox\pandoc@box{#1}%
  \Gscale@div\@tempa{\textheight}{\dimexpr\ht\pandoc@box+\dp\pandoc@box\relax}%
  \Gscale@div\@tempb{\linewidth}{\wd\pandoc@box}%
  \ifdim\@tempb\p@<\@tempa\p@\let\@tempa\@tempb\fi% select the smaller of both
  \ifdim\@tempa\p@<\p@\scalebox{\@tempa}{\usebox\pandoc@box}%
  \else\usebox{\pandoc@box}%
  \fi%
}
% Set default figure placement to htbp
\def\fps@figure{htbp}
\makeatother
\setlength{\emergencystretch}{3em} % prevent overfull lines
\providecommand{\tightlist}{%
  \setlength{\itemsep}{0pt}\setlength{\parskip}{0pt}}
\usepackage{bookmark}
\IfFileExists{xurl.sty}{\usepackage{xurl}}{} % add URL line breaks if available
\urlstyle{same}
\hypersetup{
  pdftitle={Mid-Senior Quantitative Analyst in FCP Model Validation Team \textbar{} SEB Vilnius - case study},
  pdfauthor={Kristina Kazlauskaitė},
  hidelinks,
  pdfcreator={LaTeX via pandoc}}

\title{Mid-Senior Quantitative Analyst in FCP Model Validation Team
\textbar{} SEB Vilnius - case study}
\author{Kristina Kazlauskaitė}
\date{July 2026}

\begin{document}
\maketitle

{
\setcounter{tocdepth}{2}
\tableofcontents
}
\section{1. Data understanding and
manipulation}\label{data-understanding-and-manipulation}

\subsection{1.1 Data sources and
structure}\label{data-sources-and-structure}

Two files were provided:

\texttt{alerts.xlsx} - 10,177 rows x 14 columns, one row per TM alert,
combining three logical layers: the 1st-line alert (\texttt{AlertType},
\texttt{AlertState}, \texttt{DateCreated}, \texttt{DateClosed}), the
2nd-line case investigation, populated only when an alert is escalated
(\texttt{CaseOpen}, \texttt{CaseClosed}, \texttt{CaseReported},
\texttt{CaseState}), and customer attributes at the time of the alert
(\texttt{PEP}, \texttt{CusRiskCategory}, \texttt{Type},
\texttt{IndustryCode}).

\texttt{additional\_info.xlsx} - a Metadata sheet defining each field,
and an Industry Segment sheet mapping 33 industry codes to a risk
segment (high risk / high risk cash intensive / staff intensive small
company) and a numeric risk score.

The dataset is best understood as a funnel: every row starts as an
alert; a small subset escalates to a case; a smaller subset of those
results in a SAR. The customer segment (\texttt{Type}) is the key
structural split - Private Banking (individual clients) and LC\&FI
(corporate/institutional clients) run different TM scenarios and
populate different attribute fields, so the analysis below is cut by
\texttt{Type} rather than pooled wherever it matters.

\subsection{1.2 Data quality
observations}\label{data-quality-observations}

\begin{enumerate}
\def\labelenumi{(\alph{enumi})}
\tightlist
\item
  Inconsistent missing-value encoding
\end{enumerate}

45,319 cells in the raw file contain the literal text ``NULL'' rather
than a blank cell, spread across \texttt{DateClosed}, \texttt{CaseOpen},
\texttt{CaseClosed}, \texttt{CaseReported}, \texttt{CaseState},
\texttt{PEP} and \texttt{CusRiskCategory.} Left unhandled, this would be
read as a text category or block date parsing entirely. Both blank cells
and the string ``NULL'' were mapped to true NA on import.

\begin{enumerate}
\def\labelenumi{(\alph{enumi})}
\setcounter{enumi}{1}
\tightlist
\item
  Duplicate records
\end{enumerate}

158 rows (79 distinct records) are complete duplicates of another row
once the internal export index is ignored. These were removed, leaving
10,098 unique alerts (6,267 LC\&FI / 3,831 Private Banking).

\begin{enumerate}
\def\labelenumi{(\alph{enumi})}
\setcounter{enumi}{2}
\tightlist
\item
  intID is not a reliable unique key
\end{enumerate}

Only 9,194 of 10,177 raw values of intID are unique (9.7\% collide).
Investigation showed 904 of the 968 colliding groups pair one LC\&FI
record with one Private Banking record - intID is assigned independently
within each source system and simply overlaps when the two extracts are
combined; on its own it carries no duplication risk. The remaining 64
within-segment collisions were checked individually and are fully
explained by the 79 exact-duplicate records in (b) - there is no case of
two genuinely different alerts sharing the same ID within the same
segment. Practical implication: intID should never be used as a join or
dedup key without also qualifying by \texttt{Type.}

\begin{enumerate}
\def\labelenumi{(\alph{enumi})}
\setcounter{enumi}{3}
\tightlist
\item
  \texttt{CaseClosed} is never populated for SAR-filed cases
\end{enumerate}

All 47 SAR-filed alerts in the cleaned data have a missing
\texttt{CaseClosed} timestamp - a systematic workflow gap rather than
random missingness. This made the originally intended ``case-closed to
SAR-reported'' cycle time unusable; \texttt{CaseOpen} to
\texttt{CaseReported} was used as a substitute (Table 9).

\subsection{1.3 Structural vs.~genuine
missingness}\label{structural-vs.-genuine-missingness}

Two fields are only ever populated for one customer segment and should
not be read as data-quality gaps: \texttt{PEP} is populated only for
Private Banking clients, and \texttt{IndustryCode} only for LC\&FI
clients. Table 1 reports raw completeness; elsewhere in this report
structural absence is separated from genuine gaps - e.g.~the 54 alerts
with a blank CusRiskCategory, or the 938 Private Banking alerts with no
PEP flag recorded despite PEP being applicable, which is worth raising
with the data owner.

\subsection{1.4 Assumptions}\label{assumptions}

\begin{itemize}
\item
  An alert with no \texttt{CaseOpen} date is treated as never escalated
  beyond 1st line, not as ``pending'' - reasonable given the data spans
  2008-2021 and a genuinely open investigation from that period would
  itself be an anomaly (checked for: none found).
\item
  ``Escalated'' = \texttt{CaseOpen} populated; ``SAR filed'' =
  \texttt{CaseReported} populated - the cleanest available proxies for
  2nd-line escalation and confirmed suspicious activity respectively.
\item
  Exact-duplicate records are treated as data-export artefacts and
  removed, since a genuine repeat alert would be expected to carry a new
  \texttt{DateCreated.}
\item
  The Industry Segment table is assumed static over the sample period
  (no effective-dating supplied), so it is applied as a single lookup
  rather than a time-varying attribute.
\item
  Rows closed for a stated data-quality reason (``Closed - Not
  Investigated Data Quality'', 6 LC\&FI alerts) are kept in the base
  population, since excluding them would understate true alert volume.
\end{itemize}

\begin{longtable}[]{@{}lrr@{}}
\caption{Table 1. Field completeness}\tabularnewline
\toprule\noalign{}
Field & N\_missing & Pct\_missing \\
\midrule\noalign{}
\endfirsthead
\toprule\noalign{}
Field & N\_missing & Pct\_missing \\
\midrule\noalign{}
\endhead
\bottomrule\noalign{}
\endlastfoot
intID & 0 & 0.0 \\
AlertType & 0 & 0.0 \\
AlertState & 0 & 0.0 \\
DateCreated & 0 & 0.0 \\
DateClosed & 188 & 1.9 \\
CaseOpen & 9921 & 98.2 \\
CaseClosed & 9968 & 98.7 \\
CaseReported & 10051 & 99.5 \\
CaseState & 9921 & 98.2 \\
PEP & 7205 & 71.4 \\
CusRiskCategory & 54 & 0.5 \\
Type & 0 & 0.0 \\
IndustryCode & 7907 & 78.3 \\
\end{longtable}

\emph{Completeness by field. ``Structural'' = field is not applicable to
that customer segment by design, not a data quality defect.}

\begin{longtable}[]{@{}llr@{}}
\caption{Table 2. Alert volume by customer segment}\tabularnewline
\toprule\noalign{}
Type & AlertType & n\_alerts \\
\midrule\noalign{}
\endfirsthead
\toprule\noalign{}
Type & AlertType & n\_alerts \\
\midrule\noalign{}
\endhead
\bottomrule\noalign{}
\endlastfoot
LC\&FI & Credit Cards & 1959 \\
LC\&FI & Check Countries List & 1367 \\
LC\&FI & Unusual behaviour & 843 \\
LC\&FI & Cash & 637 \\
LC\&FI & Listed High Risk Banks & 549 \\
LC\&FI & Awakening Account & 270 \\
LC\&FI & Close Monitoring & 184 \\
LC\&FI & repayement of funds & 176 \\
LC\&FI & PEP Monitoring & 168 \\
LC\&FI & Unusual Cash Behaviour & 64 \\
LC\&FI & Existing Accounts & 50 \\
Private Banking & Unusual behaviour & 1711 \\
Private Banking & New Destinations with high turnover & 1538 \\
Private Banking & Existing Accounts & 220 \\
Private Banking & Awakening Account & 159 \\
Private Banking & Recurring In-Out scenario & 97 \\
Private Banking & Check Countries List & 38 \\
Private Banking & International Transfers & 34 \\
Private Banking & PEP Monitoring & 28 \\
Private Banking & Close Monitoring & 6 \\
\end{longtable}

\begin{longtable}[]{@{}
  >{\raggedright\arraybackslash}p{(\linewidth - 12\tabcolsep) * \real{0.1600}}
  >{\raggedleft\arraybackslash}p{(\linewidth - 12\tabcolsep) * \real{0.0900}}
  >{\raggedleft\arraybackslash}p{(\linewidth - 12\tabcolsep) * \real{0.1200}}
  >{\raggedleft\arraybackslash}p{(\linewidth - 12\tabcolsep) * \real{0.1600}}
  >{\raggedleft\arraybackslash}p{(\linewidth - 12\tabcolsep) * \real{0.0600}}
  >{\raggedleft\arraybackslash}p{(\linewidth - 12\tabcolsep) * \real{0.1900}}
  >{\raggedleft\arraybackslash}p{(\linewidth - 12\tabcolsep) * \real{0.2200}}@{}}
\caption{Table 3. Alert -\textgreater{} Case -\textgreater{} SAR funnel,
by segment}\tabularnewline
\toprule\noalign{}
\begin{minipage}[b]{\linewidth}\raggedright
Type
\end{minipage} & \begin{minipage}[b]{\linewidth}\raggedleft
n\_alerts
\end{minipage} & \begin{minipage}[b]{\linewidth}\raggedleft
n\_escalated
\end{minipage} & \begin{minipage}[b]{\linewidth}\raggedleft
escalation\_rate
\end{minipage} & \begin{minipage}[b]{\linewidth}\raggedleft
n\_sar
\end{minipage} & \begin{minipage}[b]{\linewidth}\raggedleft
sar\_rate\_of\_alerts
\end{minipage} & \begin{minipage}[b]{\linewidth}\raggedleft
sar\_rate\_of\_escalated
\end{minipage} \\
\midrule\noalign{}
\endfirsthead
\toprule\noalign{}
\begin{minipage}[b]{\linewidth}\raggedright
Type
\end{minipage} & \begin{minipage}[b]{\linewidth}\raggedleft
n\_alerts
\end{minipage} & \begin{minipage}[b]{\linewidth}\raggedleft
n\_escalated
\end{minipage} & \begin{minipage}[b]{\linewidth}\raggedleft
escalation\_rate
\end{minipage} & \begin{minipage}[b]{\linewidth}\raggedleft
n\_sar
\end{minipage} & \begin{minipage}[b]{\linewidth}\raggedleft
sar\_rate\_of\_alerts
\end{minipage} & \begin{minipage}[b]{\linewidth}\raggedleft
sar\_rate\_of\_escalated
\end{minipage} \\
\midrule\noalign{}
\endhead
\bottomrule\noalign{}
\endlastfoot
LC\&FI & 6267 & 143 & 0.0228179 & 45 & 0.0071805 & 0.3146853 \\
Private Banking & 3831 & 34 & 0.0088750 & 2 & 0.0005221 & 0.0588235 \\
Overall & 10098 & 177 & 0.0175282 & 47 & 0.0046544 & 0.2655367 \\
\end{longtable}

\emph{``SAR / Escalated'' is the conditional yield of an escalation -
i.e.~once a case is opened, how often does it end in a SAR.}

\begin{longtable}[]{@{}
  >{\raggedright\arraybackslash}p{(\linewidth - 14\tabcolsep) * \real{0.1176}}
  >{\raggedright\arraybackslash}p{(\linewidth - 14\tabcolsep) * \real{0.2647}}
  >{\raggedleft\arraybackslash}p{(\linewidth - 14\tabcolsep) * \real{0.0662}}
  >{\raggedleft\arraybackslash}p{(\linewidth - 14\tabcolsep) * \real{0.0882}}
  >{\raggedleft\arraybackslash}p{(\linewidth - 14\tabcolsep) * \real{0.1176}}
  >{\raggedleft\arraybackslash}p{(\linewidth - 14\tabcolsep) * \real{0.0441}}
  >{\raggedleft\arraybackslash}p{(\linewidth - 14\tabcolsep) * \real{0.1397}}
  >{\raggedleft\arraybackslash}p{(\linewidth - 14\tabcolsep) * \real{0.1618}}@{}}
\caption{Table 4. Scenario-level productivity (Alert -\textgreater{}
Case -\textgreater{} SAR), by segment}\tabularnewline
\toprule\noalign{}
\begin{minipage}[b]{\linewidth}\raggedright
Type
\end{minipage} & \begin{minipage}[b]{\linewidth}\raggedright
AlertType
\end{minipage} & \begin{minipage}[b]{\linewidth}\raggedleft
n\_alerts
\end{minipage} & \begin{minipage}[b]{\linewidth}\raggedleft
n\_escalated
\end{minipage} & \begin{minipage}[b]{\linewidth}\raggedleft
escalation\_rate
\end{minipage} & \begin{minipage}[b]{\linewidth}\raggedleft
n\_sar
\end{minipage} & \begin{minipage}[b]{\linewidth}\raggedleft
sar\_rate\_of\_alerts
\end{minipage} & \begin{minipage}[b]{\linewidth}\raggedleft
sar\_rate\_of\_escalated
\end{minipage} \\
\midrule\noalign{}
\endfirsthead
\toprule\noalign{}
\begin{minipage}[b]{\linewidth}\raggedright
Type
\end{minipage} & \begin{minipage}[b]{\linewidth}\raggedright
AlertType
\end{minipage} & \begin{minipage}[b]{\linewidth}\raggedleft
n\_alerts
\end{minipage} & \begin{minipage}[b]{\linewidth}\raggedleft
n\_escalated
\end{minipage} & \begin{minipage}[b]{\linewidth}\raggedleft
escalation\_rate
\end{minipage} & \begin{minipage}[b]{\linewidth}\raggedleft
n\_sar
\end{minipage} & \begin{minipage}[b]{\linewidth}\raggedleft
sar\_rate\_of\_alerts
\end{minipage} & \begin{minipage}[b]{\linewidth}\raggedleft
sar\_rate\_of\_escalated
\end{minipage} \\
\midrule\noalign{}
\endhead
\bottomrule\noalign{}
\endlastfoot
LC\&FI & Credit Cards & 1959 & 78 & 0.0398162 & 27 & 0.0137825 &
0.3461538 \\
LC\&FI & Check Countries List & 1367 & 35 & 0.0256035 & 4 & 0.0029261 &
0.1142857 \\
LC\&FI & Unusual behaviour & 843 & 8 & 0.0094899 & 0 & 0.0000000 &
0.0000000 \\
LC\&FI & Cash & 637 & 14 & 0.0219780 & 11 & 0.0172684 & 0.7857143 \\
LC\&FI & Listed High Risk Banks & 549 & 3 & 0.0054645 & 3 & 0.0054645 &
1.0000000 \\
LC\&FI & Awakening Account & 270 & 1 & 0.0037037 & 0 & 0.0000000 &
0.0000000 \\
LC\&FI & Close Monitoring & 184 & 0 & 0.0000000 & 0 & 0.0000000 & NA \\
LC\&FI & repayement of funds & 176 & 2 & 0.0113636 & 0 & 0.0000000 &
0.0000000 \\
LC\&FI & PEP Monitoring & 168 & 0 & 0.0000000 & 0 & 0.0000000 & NA \\
LC\&FI & Unusual Cash Behaviour & 64 & 0 & 0.0000000 & 0 & 0.0000000 &
NA \\
LC\&FI & Existing Accounts & 50 & 2 & 0.0400000 & 0 & 0.0000000 &
0.0000000 \\
Private Banking & Unusual behaviour & 1711 & 10 & 0.0058445 & 1 &
0.0005845 & 0.1000000 \\
Private Banking & New Destinations with high turnover & 1538 & 17 &
0.0110533 & 0 & 0.0000000 & 0.0000000 \\
Private Banking & Existing Accounts & 220 & 2 & 0.0090909 & 0 &
0.0000000 & 0.0000000 \\
Private Banking & Awakening Account & 159 & 0 & 0.0000000 & 0 &
0.0000000 & NA \\
Private Banking & Recurring In-Out scenario & 97 & 0 & 0.0000000 & 0 &
0.0000000 & NA \\
Private Banking & Check Countries List & 38 & 5 & 0.1315789 & 1 &
0.0263158 & 0.2000000 \\
Private Banking & International Transfers & 34 & 0 & 0.0000000 & 0 &
0.0000000 & NA \\
Private Banking & PEP Monitoring & 28 & 0 & 0.0000000 & 0 & 0.0000000 &
NA \\
Private Banking & Close Monitoring & 6 & 0 & 0.0000000 & 0 & 0.0000000 &
NA \\
\end{longtable}

\emph{Full scenario breakdown. ``n/a'' = zero escalations, so a
conditional yield cannot be computed. See Section 3.2 for
interpretation.}

\begin{longtable}[]{@{}llrr@{}}
\caption{Table 5. 1st-line alert disposition (AlertState), by
segment}\tabularnewline
\toprule\noalign{}
Type & AlertState & n & pct \\
\midrule\noalign{}
\endfirsthead
\toprule\noalign{}
Type & AlertState & n & pct \\
\midrule\noalign{}
\endhead
\bottomrule\noalign{}
\endlastfoot
LC\&FI & Closed - Not Suspicious & 5326 & 85.0 \\
LC\&FI & Closed - Not Investigated & 471 & 7.5 \\
LC\&FI & Closed - Processed externally & 302 & 4.8 \\
LC\&FI & Data Created & 147 & 2.3 \\
LC\&FI & Unassigned & 9 & 0.1 \\
LC\&FI & Closed - Not Investigated Data Quality & 6 & 0.1 \\
LC\&FI & Assigned to Investigate & 3 & 0.0 \\
LC\&FI & Under Investigation & 2 & 0.0 \\
LC\&FI & Level 2 escalation & 1 & 0.0 \\
Private Banking & Closed - Not Suspicious & 3752 & 97.9 \\
Private Banking & Closed - Not Investigated & 38 & 1.0 \\
Private Banking & Data Created & 34 & 0.9 \\
Private Banking & Under Investigation & 4 & 0.1 \\
Private Banking & Unassigned & 3 & 0.1 \\
\end{longtable}

\emph{``Closed - Not Investigated'' and ``Closed - Processed
externally'' together account for 12.4\% of LC\&FI alerts (777 alerts)
that were closed without a substantive 1st-line review - see Section
3.3.}

\begin{longtable}[]{@{}lrrr@{}}
\caption{Table 6. Escalation / SAR rate by customer risk
category}\tabularnewline
\toprule\noalign{}
CusRiskCategory\_clean & n\_alerts & escalation\_rate & sar\_rate \\
\midrule\noalign{}
\endfirsthead
\toprule\noalign{}
CusRiskCategory\_clean & n\_alerts & escalation\_rate & sar\_rate \\
\midrule\noalign{}
\endhead
\bottomrule\noalign{}
\endlastfoot
Medium Risk & 6942 & 0.0142610 & 0.0025929 \\
Lower Risk & 1314 & 0.0197869 & 0.0038052 \\
Higher Risk & 1161 & 0.0439276 & 0.0206718 \\
Not Specified & 627 & 0.0015949 & 0.0000000 \\
Unknown / not recorded & 54 & 0.0000000 & 0.0000000 \\
\end{longtable}

\begin{longtable}[]{@{}lrrr@{}}
\caption{Table 7. Escalation / SAR rate by PEP status (Private Banking
only)}\tabularnewline
\toprule\noalign{}
PEP\_clean & n\_alerts & escalation\_rate & sar\_rate \\
\midrule\noalign{}
\endfirsthead
\toprule\noalign{}
PEP\_clean & n\_alerts & escalation\_rate & sar\_rate \\
\midrule\noalign{}
\endhead
\bottomrule\noalign{}
\endlastfoot
Not PEP & 2776 & 0.0104467 & 0.0007205 \\
PEP & 117 & 0.0170940 & 0.0000000 \\
Unknown / not recorded & 938 & 0.0031983 & 0.0000000 \\
\end{longtable}

\emph{PEP sample size (117) is too small to draw firm conclusions - see
Section 3.3.}

\begin{longtable}[]{@{}
  >{\raggedright\arraybackslash}p{(\linewidth - 6\tabcolsep) * \real{0.5270}}
  >{\raggedleft\arraybackslash}p{(\linewidth - 6\tabcolsep) * \real{0.1216}}
  >{\raggedleft\arraybackslash}p{(\linewidth - 6\tabcolsep) * \real{0.2162}}
  >{\raggedleft\arraybackslash}p{(\linewidth - 6\tabcolsep) * \real{0.1351}}@{}}
\caption{Table 8. Escalation / SAR rate by industry risk segment (LC\&FI
only)}\tabularnewline
\toprule\noalign{}
\begin{minipage}[b]{\linewidth}\raggedright
IndustrySegment\_clean
\end{minipage} & \begin{minipage}[b]{\linewidth}\raggedleft
n\_alerts
\end{minipage} & \begin{minipage}[b]{\linewidth}\raggedleft
escalation\_rate
\end{minipage} & \begin{minipage}[b]{\linewidth}\raggedleft
sar\_rate
\end{minipage} \\
\midrule\noalign{}
\endfirsthead
\toprule\noalign{}
\begin{minipage}[b]{\linewidth}\raggedright
IndustrySegment\_clean
\end{minipage} & \begin{minipage}[b]{\linewidth}\raggedleft
n\_alerts
\end{minipage} & \begin{minipage}[b]{\linewidth}\raggedleft
escalation\_rate
\end{minipage} & \begin{minipage}[b]{\linewidth}\raggedleft
sar\_rate
\end{minipage} \\
\midrule\noalign{}
\endhead
\bottomrule\noalign{}
\endlastfoot
Other / not in high-risk industry list & 6188 & 0.0224628 & 0.0072721 \\
high risk & 57 & 0.0701754 & 0.0000000 \\
staff intensive small company & 22 & 0.0000000 & 0.0000000 \\
\end{longtable}

\emph{Flagged high-risk industries escalate 3x more often but produced
zero SARs in this sample (n=57) - see Section 3.3 on sample-size
caveats.}

\begin{longtable}[]{@{}
  >{\raggedleft\arraybackslash}p{(\linewidth - 14\tabcolsep) * \real{0.1330}}
  >{\raggedleft\arraybackslash}p{(\linewidth - 14\tabcolsep) * \real{0.1232}}
  >{\raggedleft\arraybackslash}p{(\linewidth - 14\tabcolsep) * \real{0.1182}}
  >{\raggedleft\arraybackslash}p{(\linewidth - 14\tabcolsep) * \real{0.1281}}
  >{\raggedleft\arraybackslash}p{(\linewidth - 14\tabcolsep) * \real{0.1182}}
  >{\raggedright\arraybackslash}p{(\linewidth - 14\tabcolsep) * \real{0.1379}}
  >{\raggedright\arraybackslash}p{(\linewidth - 14\tabcolsep) * \real{0.1281}}
  >{\raggedleft\arraybackslash}p{(\linewidth - 14\tabcolsep) * \real{0.1133}}@{}}
\caption{Table 9. Cycle times (days)}\tabularnewline
\toprule\noalign{}
\begin{minipage}[b]{\linewidth}\raggedleft
median\_days\_to\_alert\_close
\end{minipage} & \begin{minipage}[b]{\linewidth}\raggedleft
mean\_days\_to\_alert\_close
\end{minipage} & \begin{minipage}[b]{\linewidth}\raggedleft
p90\_days\_to\_alert\_close
\end{minipage} & \begin{minipage}[b]{\linewidth}\raggedleft
median\_days\_case\_investig
\end{minipage} & \begin{minipage}[b]{\linewidth}\raggedleft
mean\_days\_case\_investig
\end{minipage} & \begin{minipage}[b]{\linewidth}\raggedright
median\_days\_close\_to\_report
\end{minipage} & \begin{minipage}[b]{\linewidth}\raggedright
mean\_days\_close\_to\_report
\end{minipage} & \begin{minipage}[b]{\linewidth}\raggedleft
n\_negative\_cycle\_times
\end{minipage} \\
\midrule\noalign{}
\endfirsthead
\toprule\noalign{}
\begin{minipage}[b]{\linewidth}\raggedleft
median\_days\_to\_alert\_close
\end{minipage} & \begin{minipage}[b]{\linewidth}\raggedleft
mean\_days\_to\_alert\_close
\end{minipage} & \begin{minipage}[b]{\linewidth}\raggedleft
p90\_days\_to\_alert\_close
\end{minipage} & \begin{minipage}[b]{\linewidth}\raggedleft
median\_days\_case\_investig
\end{minipage} & \begin{minipage}[b]{\linewidth}\raggedleft
mean\_days\_case\_investig
\end{minipage} & \begin{minipage}[b]{\linewidth}\raggedright
median\_days\_close\_to\_report
\end{minipage} & \begin{minipage}[b]{\linewidth}\raggedright
mean\_days\_close\_to\_report
\end{minipage} & \begin{minipage}[b]{\linewidth}\raggedleft
n\_negative\_cycle\_times
\end{minipage} \\
\midrule\noalign{}
\endhead
\bottomrule\noalign{}
\endlastfoot
4 & 31.7 & 103 & 0.1137635 & 237.2 & NA & NA & 0 \\
\end{longtable}

\emph{``Case opened -\textgreater{} Case closed'' mean is inflated by a
cluster of \textasciitilde10+ cases sitting at 1,153-1,155 days
(\textasciitilde3.2 years) with near-identical durations - consistent
with a backlog of aged cases being cleared in one batch rather than
measurement error. Median-based figures are more representative of
typical processing time.}

\section{2. Data Visualisation}\label{data-visualisation}

\subsection{2.1 Alert volume by scenario and over
time}\label{alert-volume-by-scenario-and-over-time}

\begin{figure}
\centering
\pandocbounded{\includegraphics[keepaspectratio,alt={Figure 1: Alert volume by TM scenario.}]{r here("figures", "F1_volume_by_scenario.png")}}
\caption{Figure 1: Alert volume by TM scenario.}
\end{figure}

Volume is highly concentrated: the top 4 scenarios (Unusual behaviour,
Credit Cards, New Destinations with high turnover, Check Countries List)
account for roughly 68\% of all alerts. The scenario mix barely overlaps
between segments - only Unusual behaviour, Existing Accounts, Check
Countries List, Awakening Account, and PEP Monitoring / Close Monitoring
run in both, and even then with very different relative volumes. This
confirms the two segments should be evaluated as separate TM models, not
pooled.

\begin{figure}
\centering
\pandocbounded{\includegraphics[keepaspectratio,alt={Figure 2: Monthly trend.}]{r here("figures", "F2_monthly_trend.png")}}
\caption{Figure 2: Monthly trend.}
\end{figure}

Volumes are volatile with several sharp, short-lived spikes (e.g.~LC\&FI
in mid-2009 and mid-2015) that look more like one-off batch events (a
rule re-run, a scenario back-test, or a system migration) than organic
growth. LC\&FI volume was consistently higher and more volatile than
Private Banking until around 2019-2020, when Private Banking volume
increased and LC\&FI's dropped - worth confirming with the business
whether this reflects a genuine change in scenario tuning, thresholds,
or client-base composition around that time.

\subsection{2.2 The alert funnel}\label{the-alert-funnel}

\begin{figure}
\centering
\pandocbounded{\includegraphics[keepaspectratio,alt={Figure 3: Alert funnel.}]{r here("figures", "F3_funnel.png")}}
\caption{Figure 3: Alert funnel.}
\end{figure}

On a log scale, each stage of the funnel drops by roughly one order of
magnitude. This shape is typical of rule-based TM - the vast majority of
alerts are closed at 1st line without escalation - and is not itself a
red flag; the question a validator should ask is whether the size of the
drop-off is appropriate for the scenario, which Figure 4 addresses
directly.

\subsection{2.3 Scenario productivity}\label{scenario-productivity}

\begin{figure}
\centering
\pandocbounded{\includegraphics[keepaspectratio,alt={Figure 4: Scenario productivity.}]{r here("figures", "F4_scenario_yield_scatter.png")}}
\caption{Figure 4: Scenario productivity.}
\end{figure}

This is the most decision-relevant chart in the report. Scenarios
cluster in the bottom-left (low escalation, low SAR rate - the
majority), with a handful of standouts: Cash and Credit Cards combine
meaningful volume with a comparatively high SAR rate; Check Countries
List (Private Banking) has by far the highest escalation rate (13.2\%)
but is low-volume (38 alerts) and its single SAR result should be read
cautiously. Scenarios sitting on the x-axis with zero height
(e.g.~Unusual behaviour LC\&FI, New Destinations with high turnover)
escalate regularly but have never converted to a SAR - see Section 3.2
for what that implies for tuning.

\subsection{2.4 Risk-based
discrimination}\label{risk-based-discrimination}

\begin{figure}
\centering
\pandocbounded{\includegraphics[keepaspectratio,alt={Figure 5: SAR rate by risk.}]{r here("figures", "F5_sar_rate_by_risk_category.png")}}
\caption{Figure 5: SAR rate by risk.}
\end{figure}

SAR rate increases with customer risk category, roughly 8x higher for
Higher Risk than for Medium Risk. This is a genuinely encouraging
validation signal: whatever combination of KYC risk scoring and TM rules
is in place, it is - on this metric at least - correctly concentrating
true positives among higher-risk customers rather than spreading them
evenly.

\subsection{2.5 First-line disposition and cycle
times}\label{first-line-disposition-and-cycle-times}

\begin{figure}
\centering
\pandocbounded{\includegraphics[keepaspectratio,alt={Figure 6: Alert disposition.}]{r here("figures", "F6_alertstate_composition.png")}}
\caption{Figure 6: Alert disposition.}
\end{figure}

Private Banking alerts are closed as ``Not Suspicious'' almost uniformly
(97.9\%). LC\&FI shows more variety, including the 12.4\% closed as
``Not Investigated'' or ``Processed externally'' flagged in Table 5 -
alerts that received no substantive first-line review, which matters for
how confidently ``closed, not suspicious'' can be read as a true
negative in any later false-negative analysis.

\begin{figure}
\centering
\pandocbounded{\includegraphics[keepaspectratio,alt={Figure 7: Time to close.}]{r here("figures", "F7_cycle_time_alert_close.png")}}
\caption{Figure 7: Time to close.}
\end{figure}

Both segments show the expected shape - most alerts closed within days -
but with a long tail extending to 100+ days, and a distinct secondary
bump around 130-140 days for LC\&FI that likely corresponds to the same
aged-case clearance pattern noted in Table 9.

\begin{figure}
\centering
\pandocbounded{\includegraphics[keepaspectratio,alt={Figure 8: Yield trend.}]{r here("figures", "F8_yield_trend.png")}}
\caption{Figure 8: Yield trend.}
\end{figure}

Neither rate shows a stable trend - both fluctuate year to year with a
low point around 2020 (possibly volume or resourcing effects) and a
recovery in 2021. Pre-2013 years were excluded as too sparse (annual
volumes under 620 alerts) to produce a meaningful rate.

\section{3. Additional analysis -- Transaction Monitoring validation
perspective}\label{additional-analysis-transaction-monitoring-validation-perspective}

\subsection{3.1 Analyses relevant to validating a rule-based TM
model}\label{analyses-relevant-to-validating-a-rule-based-tm-model}

A model validation of a rule-based TM system typically covers the
following areas. Those marked (*) are illustrated with this dataset in
Section 3.2; the rest would require data not available here
(transaction-level detail, negative/below-the-line population, threshold
parameters), and are flagged as such.

\begin{itemize}
\item
  Conceptual soundness - do the scenarios map to documented, relevant
  typologies for each customer segment, and is the logic (thresholds,
  look-back windows, aggregation rules) reasonable given the segment's
  risk profile?
\item
  Productivity / yield analysis (*) - alert volume, escalation rate, and
  SAR conversion rate, overall and by scenario. This is the core
  effectiveness metric and is covered in depth below.
\item
  Above-the-line / below-the-line (ATL/BTL) testing - sampling
  transactions that fall just below a threshold (i.e.~would not have
  alerted) to estimate the false-negative rate. Not possible here since
  only already-generated alerts are provided, not the underlying
  transaction population - see Section 3.3.
\item
  Threshold sensitivity / tuning analysis - re-running scenario logic at
  different threshold values to trace out a volume-vs-yield curve and
  identify an efficient operating point. Requires access to the
  underlying transaction data and rule logic, not available here.
\item
  Segmentation and stability analysis (*) - does performance hold
  consistently across customer segments, risk categories, PEP status,
  industry, and over time, or does it concentrate in a few
  sub-populations (which could indicate a scenario should be narrowed or
  split)?
\item
  Data quality / input integrity review (*) - are the fields feeding the
  rules complete, consistent, and correctly typed? Weaknesses here
  directly undermine rule reliability regardless of how well-designed
  the logic is.
\item
  Coverage / scenario-gap analysis - benchmarking the scenario library
  against known typologies (e.g.~FATF guidance, peer institution
  scenario libraries) to identify typologies with no corresponding rule
  at all. Requires an external typology inventory, not available here.
\item
  Timeliness / process-efficiency review (*) - cycle times from alert to
  closure and from case to reporting, since TM effectiveness includes
  acting on true positives fast enough to be useful to the FIU.
\item
  Peer / regulatory benchmarking - comparing alert-to-SAR conversion
  rates against industry peers or supervisory expectations. Not possible
  without external benchmark data.
\item
  Independent testing / quality assurance sampling - periodic
  2nd/3rd-line re-review of a sample of 1st-line ``closed, not
  suspicious'' alerts to estimate the false-negative rate within the
  alerted population itself (a cheaper partial substitute for full
  ATL/BTL testing).
\end{itemize}

\subsection{3.2 What this data says about TM scenario
performance}\label{what-this-data-says-about-tm-scenario-performance}

\begin{itemize}
\item
  Overall productivity is low but not unusual for rule-based TM. A
  1.75\% escalation rate and 0.47\% alert-to-SAR rate implies a
  false-positive rate (at 1st line) north of 98\%. This is broadly
  consistent with published industry experience for rule-based TM
  (commonly cited false-positive rates are in the 90-98\%+ range), so on
  its own this is not evidence of a poorly calibrated model - but it
  does mean the model's practical value depends heavily on whether the
  \textasciitilde2\% of alerts that do escalate are the right 2\%.
\item
  The risk-based signal is real (Table 6, Figure 5), but PEP and
  Industry results should be read cautiously. SAR rate rising
  monotonically with declared customer risk is a genuinely positive
  validation finding. The PEP result (Table 7) and high-risk-industry
  result (Table 8) point the opposite direction at first glance (PEP
  alerts and high-risk-industry alerts escalate more but yield zero
  SARs) - however, both are built on very small samples (117 and 57
  alerts respectively) and zero SARs in a small sample is not strong
  evidence of zero true underlying risk; this needs a larger sample
  before it should influence rule design.
\item
  Scenario-level performance is highly uneven (Table 4, Figure 4). Eight
  scenario/segment combinations (accounting for \textasciitilde740
  alerts, 7.3\% of volume) never escalated a single alert across the
  whole period - Close Monitoring, PEP Monitoring and Unusual Cash
  Behaviour (LC\&FI), and Awakening Account, Recurring In-Out scenario,
  International Transfers, PEP Monitoring and Close Monitoring (Private
  Banking). A further six scenarios escalate regularly but have never
  produced a SAR (Unusual behaviour and Awakening Account and repayement
  of funds and Existing Accounts LC\&FI; New Destinations with high
  turnover and Existing Accounts Private Banking - 32 escalations
  combined with zero conversions). Both groups are natural candidates
  for a threshold or logic review, discussed further in Section 3.4.
\end{itemize}

\subsection{3.3 Key limitations and
risks}\label{key-limitations-and-risks}

\begin{itemize}
\item
  No visibility into false negatives. The dataset only contains alerts
  that were already generated. There is no way to estimate how much
  suspicious activity exists in the transaction population but was never
  flagged - the single biggest limitation for any TM validation and one
  that cannot be resolved from alert data alone (see ATL/BTL testing,
  Section 3.1).
\item
  No underlying transaction or threshold data. Without the transaction
  amounts, counts, and rule parameters that generated each alert, it is
  not possible to assess threshold calibration directly or simulate
  ``what if'' tuning scenarios - only outcome rates can be observed, not
  the mechanism producing them.
\item
  Small SAR sample size. 47 SAR-filed alerts total (2 of them Private
  Banking) is a thin basis for statistically robust scenario-level
  conclusions - several of the scenario-level yield figures in Table 4
  are based on single-digit SAR counts and should be treated as
  directional, not precise.
\item
  Unreviewed alert backlog. 12.4\% of LC\&FI alerts were closed without
  a substantive 1st-line review (Table 5, Section 2.5). This weakens
  confidence that the ``closed, not suspicious'' population is genuinely
  a clean negative class - some unknown fraction of it may contain
  unreviewed true positives.
\item
  Structural missingness limits comparability. PEP is only observed for
  Private Banking and IndustryCode only for LC\&FI, so risk-driver
  analysis cannot be run consistently across both segments with the same
  variables.
\item
  Long, volatile historical span. The 2008-2021 window likely spans
  multiple scenario logic revisions, threshold changes, and possibly
  system migrations (Figure 2's spikes and step-changes are consistent
  with this). Pooling across the full period risks averaging over
  genuinely different model states; a validator would want the change
  history to interpret trends correctly.
\item
  Data quality gaps identified in Section 1.2 (inconsistent NULL
  encoding, the CaseClosed workflow gap, and the non-unique intID) would
  each need remediation at source, since a validator relying on this
  data for regulatory reporting or model performance attestation should
  not have to first reverse-engineer these issues.
\end{itemize}

\subsection{3.4 Potential improvements to the current rule-based
model}\label{potential-improvements-to-the-current-rule-based-model}

\begin{itemize}
\item
  Retire or redesign zero-yield scenarios. The eight never-escalating
  scenarios (Section 3.2) should be reviewed with the business: either
  the underlying typology risk is genuinely very low for that segment
  (in which case retire or lower priority), or the rule logic/threshold
  is set too conservatively to ever trigger (in which case retune).
  Either conclusion requires threshold/transaction data not available
  here.
\item
  Investigate scenarios that escalate but never convert. The six
  scenarios with escalations but zero SARs (32 combined escalations)
  suggest the escalation threshold may be too loose relative to what
  2nd-line investigators consider suspicious - worth a root-cause review
  of a sample of those specific cases.
\item
  Close the unreviewed-alert gap. The 12.4\% of LC\&FI alerts closed
  without substantive review (Table 5) is an operational risk as much as
  a model one - it should be resourced down, or the reason coded
  systematically (backlog vs.~genuine low-priority triage) so a
  validator can distinguish the two going forward.
\item
  Move from static thresholds toward risk-weighted or dynamic
  thresholds. Given the clear risk-based signal in Figure 5, thresholds
  calibrated per customer-risk tier (rather than one threshold per
  scenario applied uniformly) could plausibly reduce false positives for
  Medium/Lower risk customers without giving up sensitivity for Higher
  risk ones - this would need to be tested via threshold-sensitivity
  analysis once transaction data is available.
\item
  Introduce systematic ATL/BTL and QA sampling. Regular sampling just
  below alert thresholds, and periodic 2nd-line re-review of a sample of
  ``closed, not suspicious'' 1st-line decisions, would give an estimate
  of the false-negative rate that is currently completely unmeasured.
\item
  Fix the underlying data issues. Consistent NA encoding, a populated
  CaseClosed field for SAR-filed cases, and a genuinely unique alert key
  would materially reduce the effort (and risk of misinterpretation) in
  every future validation cycle.
\end{itemize}

\subsection{3.5 Conceptual challenger model design and validation
approach}\label{conceptual-challenger-model-design-and-validation-approach}

\subsubsection{Design}\label{design}

A practical challenger to a rule-based TM model does not need to be a
single model - a staged approach is both more realistic to validate and
easier to get comfortable with from a governance perspective:

\begin{itemize}
\item
  Stage 1 - supervised risk-scoring model run in parallel (shadow mode).
  Train a model (e.g.~gradient-boosted trees, or logistic regression as
  an interpretable baseline) to predict P(SAR \textbar{} alert) using
  features engineered from the same underlying transaction/customer data
  that feeds the existing rules, plus features the current rules do not
  use explicitly - e.g.~customer risk category, PEP status, industry
  segment (already shown to carry signal in Section 3.2), account
  tenure, alert recency/frequency for that customer, and network
  features (see Stage 2). This directly targets the observed weakness:
  the current rules treat scenarios independently and do not obviously
  weight by customer risk.
\item
  Stage 2 - network/graph-based analytics as a complementary signal.
  Rule-based scenarios are inherently single-customer,
  single-transaction pattern detectors and structurally cannot see
  multi-party layering or structuring across related accounts. A
  graph-based anomaly layer (e.g.~flagging unusual counterparty
  concentration, circular flows, or connections to already-SAR'd
  entities) targets a typology gap the current model cannot see by
  construction, rather than just re-scoring the same population.
\item
  Stage 3 - combine into a blended alert score / triage layer. Rather
  than replacing the rules outright (which would require re-establishing
  regulatory comfort from zero), use the challenger's score to triage
  rule-based alerts - prioritising investigator time toward alerts the
  model scores as high-likelihood, and potentially auto-closing or
  batch-reviewing the lowest-scored tail with appropriate governance
  sign-off. This keeps the rules as the auditable, explainable trigger
  while using the challenger to address the productivity problem.
\end{itemize}

\subsubsection{Validation of the challenger against the existing
model}\label{validation-of-the-challenger-against-the-existing-model}

\begin{itemize}
\item
  Back-testing on historical alerts. Score every historical alert in
  this dataset (or a hold-out period) with the challenger model and
  compare its ranking against actual outcomes (escalated / SAR-filed).
  Standard classifier metrics apply: AUROC/AUPRC (AUPRC particularly
  relevant given the \textasciitilde0.5\% base rate), precision-recall
  at plausible operating points, and calibration (does a ``70\% risk''
  bucket actually convert to SAR \textasciitilde70\% of the time).
\item
  Champion-challenger overlap analysis. For a period run in shadow mode,
  quantify: (i) the overlap between rule-flagged and challenger-flagged
  populations, (ii) SARs captured by both vs.~only-rules
  vs.~only-challenger, and (iii) whether the challenger would have
  caught any of the historical SARs that came from scenarios with
  otherwise poor conditional yield (Table 4) - this directly tests
  whether the challenger adds coverage the rules miss, not just re-ranks
  the same alerts.
\item
  False-negative estimation via sampling. Since the challenger, unlike
  the rules, can score the full transaction population (not just
  already-alerted transactions), it can be used to sample a
  below-the-line population and have investigators review it - giving a
  first real estimate of the rules' false-negative rate, which as noted
  in Section 3.3 is currently unmeasured.
\item
  Stability and drift monitoring. Track score distributions, feature
  drift, and conversion-rate-by-score-band over time (paralleling Figure
  8 in this report) before any production reliance - Figure 8 shows the
  existing rules already have volatile year-on-year yield, so the
  challenger's stability bar should be at least as high, tested on
  rolling windows rather than a single train/test split.
\item
  Governance and explainability sign-off. Given this is a
  regulatory-facing control, model risk management would require:
  documented feature rationale, sign-off from AML/compliance on
  interpretability of the score (particularly if used to suppress or
  auto-close alerts), fair-lending/bias testing across customer segments
  (especially given PEP and risk-category are candidate features), and a
  defined fallback if the challenger's data pipeline fails - the
  existing rules would remain the control of record until the challenger
  has a sufficient validated track record to justify a formal
  champion-challenger swap.
\end{itemize}

\end{document}
